%%%%%
%
\section{Moyenne-Echantillonnage 
 en tol�rant des valeurs manquantes (param\`etre \texttt{float})}

\itemstyle{Op\'erateur~:} \operfmt{MoyEch\_mnq}
\label{OP-MoyEch_mnq}
\index{MoyEch\_mnq}

%%%%%OBJET%%%%
\itemstyle{Objet~:}
Moyenne glissante et \'echantillonnage d'un
param\`etre \`a une fr\'equence donn\'ee, diviseur
de la fr\'equence du param\`etre d'entr\'ee.
Le calcul tient compte de la pr�sence des valeurs manquantes
et permet � l'utilisateur de choisir la tol�rance en terme 
de nombre de valeurs manquantes.

%%%%%ARGUMENTS%%%%
\bigskip
%%%%%ENTREES%%%%
\begin{tabarg}
\multicolumn{3}{l}{\titretab {Entr\'ees}}\\
\hline
$X$ & Vecteur, \texttt{float} & Param\`etre d'entr\'ee\\
$R_f$ & Coeff., \texttt{double} & Rapport des fr\'equences $ \frac{F_{sortie}}
{F_{entree}}$, dans l'intervalle $]0, 1]$\\
$Decal$ & Coeff., \texttt{int} & D\'ecalage (\textit{nombre d'�chantillons}),
 dans l'intervalle $[-1/R_f, 1/R_f]$ \\
$Toler\_{mnq}$ & Coeff., \texttt{int} & Tol\'erance pour les valeurs manquantes
 (\textit{nombre d'�chantillons})\\
\end{tabarg}

\medskip
%%%%%SORTIES%%%%
\begin{tabarg}
\multicolumn{3}{l}{\titretab {Sorties}}\\
\hline
$X_m$ & Vecteur, \texttt{float} & Moyennes (m\^eme unit\'e que $X$)\\
\end{tabarg}


%%%%%FORMULE%%%%
\itemstyle{Formule}

%%Moy
\[
Xm[i] = R_f . \sum_{j = i/R_f + Decal}^{(i+1)/R_f -1 + Decal} X[j]
\qquad i=0,Nbre\_Echant(X). R_f
\]


%%%%COMMENTAIRE%%%%
\itemstyle{Commentaire~:}

$R_f$ doit \^etre l'inverse d'un entier. Si ce n'est pas le cas, utiliser
\hyperref[OP-MoyEch_ff]{\operfmt{MoyEch\_ff}}.  \newline
Le nombre d'\'echantillons du vecteur de sortie est le produit du
nombre d'\'echantillons du vecteur d'entr\'ee par $R_f$, arrondi
\`a l'entier sup\'erieur. \newline

%%%%%AUTEUR%%%%
\itemstyle{Auteur~:}
CNRM/GMEI/TRAMM

%%%%%REFERENCE%%%%
\itemstyle{R\'ef\'erences~:}
